\vssub
\subsection{~水深随时间变化} \label{sub:num_depth}
\vssub

水深的时间变化会导致局地波数网格改变。由于波数谱相对于水深时间变化是不变的，这对应于把谱从旧网格简单插值到新网格，而不改变谱形。如上文所述，新网格直接由全局不变的频率网格、新水深 $d$ 和色散关系式 (\ref{eq:disp}) 得到。更新水位的时间步长通常由水位变化的物理时间尺度决定，而不是由数值考虑决定 \citep{tol:GAOS98b}。

插值到新波数网格时使用一种简单的守恒插值方法。在该插值中，先将旧谱乘以谱箱宽度，转换为离散作用量密度。随后按常规线性插值方式，把该离散作用量重新分配到新网格上。新的离散作用量再除以（新的）谱箱宽度，转换回谱。由于旧网格通常不会完全覆盖新网格，该转换需要在高频和低频处对原始谱作参数化延拓。旧谱中位于低波数、且不被新网格解析的能量/作用量会被直接移除。新网格中位于低波数、且不被旧网格解析的部分假定能量/作用量为零。在新网格的高波数处，必要时采用通常的参数化谱尾。后一种修正较少发生，因为最高波数通常对应深水。

在实际应用中，网格修正通常只与少部分网格点有关。为避免不必要计算，仅当离散谱中最小相对水深 $kd$ 小于 4 时才变换网格。此外，仅当空间网格点未被冰覆盖，且最大波数变化至少为 $0.05 \Delta k$ 时，才对谱进行插值。
